% Assignment 2 report: Chapters 1--2 only.
% Insert into the existing Overleaf document after \begin{document}.
% Required packages: amsmath, amssymb, booktabs, graphicx.
% Upload the existing assignment2/figures/workspace_design.png alongside this file.
% Sources checked: assignment brief; notebook scene/acquisition cells;
% src/robot.py; src/simulation/constants.py; src/student/workflow.py;
% REPORT_NOTES.md; results/comprehensive_v1_full_summary.json.

\section{Introduction}
\label{sec:introduction}

Accurate robot positioning requires a kinematic model that accounts for
departures from nominal geometry and encoder readings. This assignment
calibrates a six-joint Cartesian robot comprising three prismatic joints
and a three-axis spherical wrist, using an optical tracker and an
end-effector marker body. The specified accuracy targets are
$0.2\,\mathrm{mm}$ in position and $10^{-3}\,\mathrm{rad}$ in orientation
throughout the robot's work volume, with efficient calibration data
collection. Both forward-pose prediction and physically executed inverse
kinematics must be evaluated.

We follow the assignment's recommended parameter-estimation approach
(Method~1). Homing measurements establish joint offsets, and individual
joint sweeps provide observations for a corrected geometric model with
smooth joint-distortion terms. The final estimator uses physically
constrained nonlinear least squares, while retaining proper rigid-body
transformations and an explicit convention for an unobservable parameter
direction. This approach avoids collecting a dense six-dimensional
calibration grid.

The evaluation separates tracker prediction, physical forward accuracy,
numerical forward/inverse consistency, and executed inverse accuracy.
Robot/tracker seeds 1, 7, and 61 are calibrated independently; ordinary
workspace tests and boundary, joint-limit, and wrist-singularity tests
are reported separately. Hidden simulator quantities are reserved for
evaluation and never enter calibration. The TA-provided inverse solver
is retained, so its numerical limitations are distinguished from the
accuracy of the calibrated forward model.

\section{System Design and Calibration Setup}
\label{sec:system-design}

\subsection{Robot model and coordinate conventions}

Let $W$, $T$, $B$, and $E$ denote the workspace, tracker, robot-base, and
end-effector frames. A transformation ${}^{A}T_{B}$ maps coordinates from
frame $B$ into frame $A$. The nominal base-to-end-effector transformation is
\begin{equation}
 F_{\mathrm{ee}}(\boldsymbol{\theta})
 =F_c(\theta_1,\theta_2,\theta_3)\,F_{cw}\,
   R_x(\theta_4)R_y(\theta_5)R_z(\theta_6)\,F_{we},
 \label{eq:nominal-chain}
\end{equation}
where $F_c=\operatorname{Trans}(\theta_1,\theta_2,\theta_3)$,
$F_{cw}=\operatorname{Trans}(500,0,0)$, and
$F_{we}=\operatorname{Trans}(100,0,0)$, with translations in millimeters.
Rotations in Eq.~\eqref{eq:nominal-chain} are understood as homogeneous
transformations. Thus, the wrist center has a fixed $500\,\mathrm{mm}$
offset along base $x$, and the end-effector origin is
$100\,\mathrm{mm}$ from the wrist center.

\begin{table}[htbp]
 \centering
 \caption{Robot joint ranges and maximum speeds used in calibration.}
 \label{tab:joint-specifications}
 \begin{tabular}{lll}
  \toprule
  Joints & Permitted range & Maximum speed \\
  \midrule
  $1,3$ & $[0,600]\,\mathrm{mm}$ & $50\,\mathrm{mm/s}$ \\
  $2$ & $[-300,300]\,\mathrm{mm}$ & $50\,\mathrm{mm/s}$ \\
  $4,5$ & $[-3\pi/4,3\pi/4]\,\mathrm{rad}$ & $0.1\,\mathrm{rad/s}$ \\
  $6$ & $[-\pi,\pi]\,\mathrm{rad}$ & $0.1\,\mathrm{rad/s}$ \\
  \bottomrule
 \end{tabular}
\end{table}

The simulator includes base-placement error, axis misalignment,
wrist-to-tool error, and systematic and random joint-coordinate errors.
Marker observations additionally contain random and spatially varying
systematic tracker errors. Camera-pose jiggle is disabled by the
assignment; marker measurements remain noisy, with per-coordinate
standard deviation $0.1\,\mathrm{mm}$. Each move incurs a
$0.05\,\mathrm{s}$ settling time in addition to its acceleration-limited
travel time, and each tracker snapshot costs $0.03\,\mathrm{s}$.
These acquisition costs are recorded separately from wall-clock
computation time and the assignment's ten-minute program-execution limit.

\subsection{Calibration marker-body geometry}

The special calibration body carries four markers in a regular
tetrahedral arrangement, expressed in frame $E$:
\begin{equation}
 \{\mathbf a_k\}_{k=1}^{4}
 =s\{(1,1,1),(1,-1,-1),(-1,1,-1),(-1,-1,1)\},
 \qquad s=37.5\,\mathrm{mm}.
 \label{eq:marker-layout}
\end{equation}
The centroid is the end-effector origin, and every marker lies
$\sqrt{3}s=64.95\,\mathrm{mm}$ from it, satisfying the
$75\,\mathrm{mm}$ radius constraint. The noncoplanar, symmetric layout
provides spatial separation for rigid registration without favoring
one rotational direction. Under the special-body API, marker coordinates
are exact and the body frame coincides with $E$. Consequently, no
marker-manufacturing refinement or additional marker-to-tool transform
is estimated.

\subsection{Tracker placement and workspace visibility}

The tracker frame is chosen to coincide with the workspace frame and
looks along $+z$. The approximate robot-base pose is
\begin{equation}
 {}^{W}T_T=I,
 \qquad {}^{W}T_{B,\mathrm{approx}}
 =\operatorname{Trans}(-800,0,1300),
 \label{eq:scene-placement}
\end{equation}
where the translation coordinates are in millimeters.
This centers nominal wrist-center travel at $(0,0,1600)\,\mathrm{mm}$,
with $x,y\in[-300,300]\,\mathrm{mm}$ and
$z\in[1300,1900]\,\mathrm{mm}$.

For any wrist orientation, the tool offset and marker radius give a
conservative wrist-center-to-marker distance
$\rho=100+64.95=164.95\,\mathrm{mm}$. Therefore, all nominal markers lie
within $|x|,|y|\leq464.95\,\mathrm{mm}$ and
$1135.05\leq z\leq2064.95\,\mathrm{mm}$.
The tracker viewing frustum spans $1000\leq z\leq2200\,\mathrm{mm}$,
with half-width
\begin{equation}
 h(z)=500+\frac{500}{1200}(z-1000)\quad\mathrm{mm},
 \label{eq:tracker-frustum}
\end{equation}
where $z$ is expressed in millimeters. At the nearest point of the marker
envelope, $h=556.27\,\mathrm{mm}$, leaving at least
$91.32\,\mathrm{mm}$ lateral clearance; near/far depth clearances are
approximately $135.05\,\mathrm{mm}$. This establishes nominal visibility
for the full joint range, rather than only at the sampled poses.
It is a geometric design bound, not a guarantee under every unknown
physical perturbation. Numerical geometry checks also retain all markers
for 81/81 initial visualization poses and a separate 3375-pose nominal sweep.

\begin{figure}[htbp]
 \centering
 \includegraphics[width=\linewidth]{workspace_design.png}
 \caption{Nominal tracker/robot placement and conservative marker envelope
 (left), and the tetrahedral calibration body in end-effector coordinates
 (right). Visibility shown here is a geometry check, not calibrated
 positioning accuracy.}
 \label{fig:calibration-setup}
\end{figure}

\subsection{Measurement acquisition and calibration initialization}

Five physical homings provide an additive offset estimate
$\widehat{\mathbf o}=-\frac{1}{5}\sum_{i=1}^{5}
\boldsymbol{\theta}_{\mathrm{home},i}$, with the minus sign following
from zero physical joint coordinates at home. Ten mixed baseline poses
provide an initial diagnostic of nominal-model error. The main calibration
then sweeps one joint at a time while holding the others at zero,
returning to the zero-encoder reference after each sweep.
The linear axes use 17 levels each ($37.5\,\mathrm{mm}$ spacing);
the wrist axes use 7, 7, and 9 levels, with angular spacing no greater
than $\pi/4$. Shared reference poses yield 69 distinct configurations.
Eleven re-commanded readings per configuration, plus six return-reference
readings, give $69\times11+6=765$ calibration observations.
Re-commanding renews robot repeatability noise as well as tracker noise;
it does not eliminate systematic errors.

Each valid tracker frame is paired with a copy of the settled public
encoder reading. Missing observations are omitted without silently
reusing an earlier frame, and all attempted motions/samples contribute
to simulated elapsed time. An initial base-to-tracker estimate averages
the transforms $F_{\mathrm{obs},j}F_{\mathrm{nom}}(\boldsymbol{\theta}_j)^{-1}$,
using a proper-rotation frame mean. The same setup and acquisition design
are used independently for each robot/tracker seed, taking approximately
$554.03$--$554.06\,\mathrm{s}$ simulated time including homing and baseline
measurements. Model fitting and its mathematical justification are
developed in the following sections.
